
\subsection{The Augsburg Understanding of the Congregation}

Again we read: “My standpoint is this. First it is necessary to become clear as to what the Augsburg understanding of the congregation is, and afterward there may be occasion to test it in the light of God’s Word.” Folkebladet, March 27.

How, then, is one to arrive at this clarity? Well, see here: “I do not hesitate to say that I cannot see how anyone can rightly evaluate Sverdrup’s ecclesiastical views unless he acquaints himself with the constitutional literature of the older Congregationalists and with Sohm’s works on church law.” Folkebladet, April 10.

That our people possessed knowledge of God’s Word, the catechism, and the Lutheran Confession, and that they evaluated “Sverdrup’s ecclesiastical views” precisely because of his unshakable adherence to these things, this he cannot see.

His insight is to serve as proof, and presumably means that when he cannot see it, then it is hidden from everyone else; but it may also be an open declaration that his sight is poor. Facts are not changed because a man declares that he cannot see them. And you Free Church man and Augsburger who have not acquainted yourself “with the constitutional literature of the older Congregationalists and with Sohm’s works on church law,” you have not “rightly evaluated Sverdrup’s ecclesiastical views,” according to Evjen’s opinion, for he cannot see how you could have done so. It is by Congregationalist polity literature that Augsburg’s understanding of the congregation is to be tested; the testing “in the light of God’s Word,” after all, never came. But this is in full agreement with what he says at the end of his essay: “This, however, I wish to emphasize: the Free Church’s Guiding Principles place the chief weight upon the congregation and its spirit, life, and freedom.”

This is without doubt to strike the right note, provided the note has not already been struck; but that is ordinarily done at the beginning, for after one has sung a long song in another key and then at the end strikes the right note, it does not tend toward edification.

\subsection{A Scholarly Conscience}

Prof. Ebjern says:

“And when I was finished with Augsburg” (it should presumably be at Augsburg; for if he was finished with Augsburg, he would hardly have returned to it so soon) “I had to continue my schooling for three or four years before I could, with a reasonably good scholarly conscience, present myself as a candidate for a pastoral call or a teaching position.”

There is surely no pastor or layman within the Lutheran Free Church who has anything to object to in a man continuing his schooling after he has finished at Augsburg, whether for three or even six years. He would not become so tremendously learned after all in comparison with what still remained to be learned. It is indeed precisely one of the principles for pastoral education at Augsburg that “everything belongs here that can be a help toward penetrating into God’s Word, and everything is excluded that can only dull heart and mind toward the Lord’s living truth. Everything is useful that can help toward insight into the Lord’s Word and his plan of salvation; everything is harmful that does not contribute thereto.” But advertising one’s learning and scholarship does not belong among Augsburg’s principles. Nor is there scarcely any scriptural condition for entering the work that one should possess “a good scholarly conscience.” A better condition is surely that the candidate know his own unfitness. “The true fitness is to know one’s own unfitness,” and in the consciousness of one’s own unfitness to seek God’s power against sin and the devil. What, indeed, will Satan care about a good scholarly conscience, if it means that the man is “reasonably” well satisfied with his learning and scholarship?

But it is not easy to know what Prof. Evjen puts into the expression. We proceed, however, from the assumption that conscience is the sore point in a human being—“the wound left by the lost image of God”—the tender place with which the Lord puts himself into relation, for which reason it has also been called “God’s tribunal in the heart.” John 8:9. Paul is the one who speaks most about conscience. It lies outside my purpose to demonstrate all the places in Scripture that speak of conscience, although I have pondered them all. But this much is certain: it is no infallible guide. Therefore it has been compared with the sundial. As the sundial shows the time only when it is illuminated by the sun, so conscience shows the right way only when it is illuminated by the Holy Spirit through the Word. I must set my clock by the sun—not the reverse.

In other words: the truth of God’s Word is not to be corrected according to custom and usage, tradition, or polity literature, and so forth, but all such things must be corrected according to the decisive authority of God’s Word. Conscience likewise. The smarting wounds need healing. But a candidate who goes forth into battle against sin and the devil and is “reasonably” well satisfied that he has found healing by the scholarly road for the smarting wounds of conscience, he is unfit to lead sinners to the cross of Christ. He is bound in his self-confidence and cannot himself go there. He would gladly become so great that he hides the cross of Christ from the little, groping, wandering children of men.

Or is it not precisely upon this point that Augsburg has fought so long and stubborn a battle and has won for itself both friends and enemies? Pride, ambition, and self-righteousness are indeed among the best things we human beings possess according to our nature. Therefore these things stand at a high price among men. But Jesus says: “That which is highly esteemed among men is an abomination before God.”

And why? Precisely because it is so high and great that it hides the cross of Christ and closes the way of faith against itself and against others. Prof. Evjen’s aim, then, was to attain a good scholarly conscience through his schooling; but at Augsburg this could not be done. He had to continue for another three or four years, and even then he became only “reasonably” so, which was, after all, quite understandable.

According to our catechetical instruction, sin has “destroyed the joy and peace of conscience,” for which reason Adam and Eve hid themselves behind the trees of the garden, pale with fear and red with shame. But in regeneration a new spiritual life is created, and man again receives true peace and joy in the conscience. This is the work of God’s Spirit and no work of man, so that no one may boast of it at the expense of others. See Folkebladet, page 252.

The annual meetings of the Lutheran Free Church are, as a rule, vigorous, awakening, and instructive meetings. Believing men and women come together voluntarily in order, through reports and accounts, to be reminded of the greatness and necessity of the work, the seriousness of the call, and their own unfitness, and together to pray for renewed courage and power, to encourage and urge one another to take hold in faith and love and to expect every good thing from the Lord; finally, to deliberate upon how the work may be done better in the coming year. Such have our annual meetings been. We admit that they are not as good as they ought to have been. There are deficiencies that ought to be remedied, and there are things about them that might with some justice be criticized, provided that the one who criticizes has done his part to improve them.

Let us hear Prof. Evjen’s description of our annual meetings:

“Now they are a kind of adding machine for the politics carried on between the annual meetings, an emergency brake which the train crew seizes every time they see ghosts on the track.”

And why are they now so poor?

Why, because there is so little understanding of the principles.

“Had the understanding of the principles been greater, the meeting at Fargo would not have witnessed that concerted, whining opposition to Birkeland and me,” says Evjen.

Here one might almost suppose that the “good scholarly conscience” should cease to be good. Or does Evjen really mean to tell us that he knows how things would have gone if they had not gone as they did? If he means that, then it is surely time that he receive the necessary information that he does not know it. There is no human being in the world who knows how things would have been if they had not been as they are. People who make no claim to being scholarly know that much.

Had the understanding of the principles been better, things could have gone quite differently; that we concede. But how, no one knows. We may guess that they could have gone in a manner that would have been more surprising to Evjen. But guesses are not scholarship.

But when Prof. Evjen has placed such great weight upon the principles and sought to show that they are constitutional principles, how great an importance their proper understanding has, and what misfortune a defective understanding of them has brought upon our annual meetings, then it has a “confusing” effect when, at the conclusion of his learned essay, he says:

“The great question is not: Is one in harmony with this or that group of principles ... but is one in accord with the future of the kingdom of God?”

This is indeed said of the matter in general; but it is the “Guiding Principles” that were under discussion here and were to be interpreted. They must therefore surely be included. Well then, the principles concern the congregation, and the congregation has been “the great question” among us. Prof. Evjen has, by his participation, labored zealously that the principles concerning the congregation should become known and understood. He has even, in the opinion of some, laid far too much weight upon this and, judging by the amount of column space, made this matter “the great question.”

But now he kills the whole thing and sets being in harmony with the principles in opposition to being in accord with the future of the kingdom of God.

The question then arises: Is it not just as great a question to be in accord with the past of the kingdom of God, from which the light of revelation streams through the New Testament? Can we be in accord with the future of the kingdom of God without being in accord with its past? And should not Prof. Evjen first have investigated whether the principles stood in harmony with the past and future of the kingdom of God before he made such a commotion?

“Be sure, that you are right, then go ahead.”

Why, then, has he fought? The strength of the principles is that they are built upon God’s Word, not upon polity literature. When therefore many have shown themselves unwilling to go the road Evjen points out—the road through the shifting labyrinths of polity literature, where, after hearing the most varied roads recommended, one stands more perplexed than before—it is because they know a safer road.

And here lies the real reason for that disinclination to submit to Evjen’s interpretation against which he has come up. This does not mean that we are so narrow-minded, so closed against light and guidance, that we are hopeless, although Prof. Evjen seems to think so. Therefore he says:

“For even among us there are those who read the principles and interpret them as a Jew with Moses’ veil reads his Old Testament, utterly certain that they possess the right understanding, see the red thread, and need no guidance.” Folkebladet, April 10.

But evidently he is not satisfied with this testimonial. He must make it still somewhat stronger in the next issue of Folkebladet, April 17:

“Where prejudice has anchored itself, where the sense of the trail is dull, indifference great, and the assertion ‘we have Abraham as our father’ powerful—there the prospects of a fruitful occupation with the principles and the polity literature are by no means promising.”

It is almost hopeless, then. So we may well deserve the professor’s “undivided sympathy and heartfelt wish for improvement,” as he expresses it elsewhere.

Of all this I will, in all modesty, say only this: if we were all learned men, and Prof. Evjen among a company of learned men could win honor and praise for his learning, then there would be something to speak of; but among a company of ignoramuses, not so much is required. Among the blind, the one-eyed man is king!

\subsection{The Untrained Layman’s Judgment}

Prof. Evjen claims that he has “maintained the religious equality of the laity with the pastors.” But as soon as the editors of Folkebladet wish to hear from the laity, Evjen closes the way with “right” and “call” and asks in complete surprise:

“But is the gift of grace which enables one to interpret the Guiding Principles the possession of all?”

Answer: No! One can best discover that by placing Evjen’s and Sverdrup’s interpretations side by side.

No gift of grace is required to make a matter tangled and striking. Any little boy can tangle a skein of thread; but when mother comes, she finds the thread by which it may be unraveled.

And it is surely a kind of mockery when he exclaims:

“And the poor learned works, which must yield place to the untrained layman’s judgment.”

The question now is: are the principles not a religious subject, or does the professor not practice what he preaches? A third alternative is excluded.

Prof. Sverdrup wrote:

“The light of the Word shines just as much for the simple as for the wise and understanding, and preferably for the former, so that no one is exempt from occupying himself seriously and honestly with all that belongs to the kingdom of God.”

These words too Evjen has neglected to emphasize. They presumably did not fit. He had the polity literature at his disposal, and thereby he possessed the condition for being able to understand and approximately the exclusive right to interpret the principles and represent Augsburg and the new tendency in questions concerning the congregation, lay activity, and ordination. There was no room for the untrained layman’s judgment.

Evjen may object: “The third alternative is not excluded. It is the gift of grace for interpreting the principles, and this is not the possession of all.”

Is there not here a confusion of gifts of grace with opportunity?

Answer: The untrained layman’s judgment has things, as a rule, more right in the heart than in the head. It feels when something merely dazzles the spirit without warming the heart. Among our laity there are quite a number who possess the gift of grace to discern spirits. And precisely because it is a gift of grace, it is not something learned or schooled.

This gift of grace enables one to close the heart when something profane seeks to force its way into the sanctuary, as the flower closes on a cool evening so that the cold shall not penetrate. Some possess more, others less, of this gift. Still others lack it and seem to believe that everything is equally good so long as it sounds pious or bears the scholarly stamp. Others again recognize nothing good outside their own party.

It is given as a free gift for protection against errors and not for the praise of its possessor. Never should any gift of grace be given for the possessor’s own glorification, so that he might shine by it and seize as spoil what was given for the benefit and edification of the congregation. If anyone wishes to think more highly of himself than he ought to think by overvaluing his gift of grace, he shall at once be admonished to walk in humility and not cause offense by making himself great at the expense of others.

If anyone is tempted to undervalue his gift, either because it was so small in the eyes of men or because its use did not win the recognition that the use of other gifts received, he shall at once be admonished to be faithful in the use of what he has, since judgment will fall not according to the greatness of the gift, but according to the possessor’s faithfulness in its use.

Therefore this gift of grace stands alongside the other gifts of grace, equally entitled and equally useful, since it prevents misuse of freedom and forms a protection against false and seductive spirits. Often these are clothed in the garb of truth and zeal and are therefore very difficult to recognize before “their fruits” become visible to all. But then it may be too late to remedy the damage. Here it is that this gift of grace is of such great importance.

It is felt as a protest against all that is false and spurious, a certain revulsion, as when one is about to eat a worm-eaten apple or sees an unpleasant insect gnawing at the heart-leaf of a rose. If a matter is somewhat unclear at the beginning, one then goes to “the fire” in order to obtain light and clarity.

“If you remain in my word, you are truly my disciples, and you shall know the truth, and the truth shall make you free.”

Therefore the point is to remain there and direct attention there if hearts are to be opened; for otherwise we shall soon see a generation growing up with large heads and small hearts, people who have something of that knowledge which puffs up, but nothing of that love which builds up.

“Two things my people have done. They have forsaken the fountain of living water and dug themselves wells—broken wells that hold no water.”

\subsection{Life and Death!}

Prof. Evjen writes in Folkebladet for March 20, 1918:

“Now there are presumably those who will call this six-foot-long case a coffin and its contents a corpse.”

Answer: Here he guesses again. He guesses what people may think of calling his bookcase and its contents. Let whoever wishes call this scholarship! But it is guesswork.

People may think of many things. They may think of saying: “No, the bookcase is no coffin; but after the doctor has performed his operation upon the ‘contents,’ there may perhaps be a demand for a coffin.”

In other words: “Quotations torn out of their connection with what precedes and follows can mutilate the whole.”

Such things people may think of saying.

But after having guessed, he takes the result of his guess and upon this supposed foundation builds the following handsome edifice of thought:

“But the strange thing about this content is that it is supposed to be a corpse in this bookcase, but immediately becomes spirit and life when it comes between the covers of Collected Writings.”

Allow me, in all modesty, to say: this matter is too serious for fooling. Therefore, joking aside.

Does Evjen mean to say that the “content” is the same, and that it has merely come “between the covers of Collected Writings”?

That becomes a rather gross accusation; and those who have taken part in publishing and selling Collected Writings have made themselves accomplices in a great literary theft, not to mention what an accusation it would be against the one who wrote it. And it would be too gross if Evjen were to mean what he says here.

What, then, does he mean?

So far as it is possible to understand him, he wishes to emphasize as strongly as he can the senselessness of setting Sverdrup’s writings so high that they “immediately become spirit and life,” while the polity literature upon the shelf is like corpses in a coffin. But he presents this in so flippant a manner that he comes to say what he “does not at all mean to say.”

Should it not be possible to write edifyingly about the congregation without such things?

The matter has its practical side. It is “strange,” and yet not altogether unknown among us. A man may stand with the Bible and attempt to expound it; but it does not take hold. It is so dry and intellectual that a believing person listening to it would rather take his New Testament and read it directly before the lecturer’s eyes than lend an ear to such things.

Another may take the Bible and deal even with the same text; but immediately it is felt that he stands in connection with the heart-life and thought-life of his hearers. These are familiar matters—and as it is with speech, so it is also with writing.

When Luther came to the University of Wittenberg, he wished at once to give the students “the kernel of the nut and the marrow of the wheat and of the bones.” Rector Pollick of Melrichstadt said of him:

“This monk will confuse all the doctors and bring forth a new doctrine and reform the whole Roman Church; for he applies himself to the writings of the prophets and apostles and builds upon the words of Jesus Christ. No one can overthrow or cut that down, neither with philosophy nor with sophistry.”

When the springs of life are opened, the thirsty may draw water with joy from the springs of salvation. There is light and life. There is power and peace and joy to be found.

Therefore there is a difference between opening the springs of life so that the light of the Gospel and the life of God may become accessible, and stopping them up by leading souls to the turbid waters of human principles.

Or, as someone has said: “It is not enough to scrape the outside of the text; we must try to get inside it.”

We do not despise the letter. It is of great use, the same use as the vessel in which I carry the water forward. Without such a vessel there would be nothing to hold the water. But the point is to penetrate through the letter and come into the spirit.

Neither do we despise form. It is necessary; but the great question will always be whether there is also that content of which the form is the expression, or whether it is merely form without content, shell without kernel, appearance without power.

Sverdrup says:

“The congregation is the assembly of the converted and believing. In this, life is the determining thing; form is accidental, in many respects entirely inessential. If anything is biblical, that is.”

Collected Writings II, 31.

We read, after all, of a congregation which had a name that it lived, but it was dead. It had certainly preserved the outward form of a congregation, just as a corpse may retain for a time the outward shape of the body.

On the other hand, a body may be deprived of both arms and legs, may be so maltreated that it has lost much of its original shape, and yet there may be life. This shows with all desirable clarity that form is not the essential thing.

But it is a familiar fact that where there is lack of inward life, there so much greater weight is laid upon outward form, in order at least to have something. And for a time it is possible to cultivate the form at the expense of the content, to lay an improper weight upon the outward and visible, so that the inward and invisible does not receive its due. This leads to death, because the often frail life did not receive the nourishment and the freedom it needed in order to grow and develop.

The purpose of form is to guard and protect life. But if that is the purpose, then it must always be adapted according to the demands of life. Life demands that the form be sufficiently spacious that it may move freely.

Yet even within the most spacious form life may sink down and die if it does not receive good and sufficient nourishment. If we are to attain the goal of life, then we must continually drink from the springs of life.

Now it is the Holy Spirit who both creates life and forms a shape for the life that has been created. Therefore we are very careful not only about life, but also about finding the “right form,” and that is the living and liberated congregation.

Our old teachers laid the weight upon spiritual life, the life that is hidden with Christ in God. Only in the power of this life does it become possible for the congregation to appear visibly in the world and in word and deed proclaim the excellencies of Christ. Therefore the congregation is life from the dead.

\subsection{The Revival}

Nearly all the congregations that joined in the Free Church work were more or less touched by the revival of the 1890s. There were birth pangs, movements, and conflicts among our people as never before. It was a wrestling struggle between life and death, between dead formalism and living faith. The brothers who stood foremost in the ranks, fought hardest, and suffered most know this.

In this struggle against death and formalism, some went too far. It appeared that they believed that if only they could have the form changed, then they would be helped. And in their striving to reject all old forms and adopt new ones, they became more occupied with forms than with anything else.

Party spirit and censoriousness entered in, and when they then found that this was not the way of the Lutheran Free Church, they broke with it and accused us of having the principles on paper but not practicing them. We wanted only to talk about the congregation, but not to have it on the prairie.

Then Prof. Evjen comes and tells us that now we shall learn why we are here and why we are fighting.

We fought for constitution. The whole thing was a constitutional difference.

But then we were not entirely alone. There were some over in Europe who fought for the same matter.

We remember the meeting at Valley City, when Prof. Evjen came, rather puffed up by his knowledge, to instruct us concerning this matter. We also remember Prof. Oftedal’s characteristic answer, part of which is reproduced here from memory:

“What that man there has now said, we went through thirty-five years ago and threw it all”—with a motion of the hand—“overboard as useless, and then we turned to this book here”—and he struck his New Testament—“in order to find what we could use.”

\subsection{Visible and Invisible}

Prof. Evjen writes in Folkebladet for April 17:

“Some of the most unfortunate expressions theologians have been accustomed to use concerning the congregation are visible and invisible, inward and outward. They have an altogether confusing effect. S. and O. have to some degree accepted this usage, and it was hardly to be expected that they could use these expressions for thirty-five years without committing one inconsistency or another—for whom have these expressions not caused to stumble?”

Answer: Evjen seems to concede that “the congregation itself has both an inward side and an outward side.” Why then should the expressions “inward and outward” be among the most unfortunate expressions?

If only they had been unfortunate, or more unfortunate; but to use the superlative in judging expressions that he himself continually employs can scarcely be recognized as scholarship by any “faculty of standing.”

He ought to avert the misfortune by putting a stop to the expressions. But by making diligent use of them, new misfortunes arise, and the “confusion” continues.

First Sverdrup and Oftedal must come before the bar as accused of linguistic abuses and inconsistencies. But the evidence against them is so inadequate that the case is dismissed. Should it be taken up again, there will be time to demonstrate this more closely.

Our time is serious, and there are other things to do than quarrel about words and expressions.

The expressions have been misused, but for that reason we cannot reject their proper use.

Different meanings may be put into them, as when the Catholics said:

“The Church is visible like the kingdom of France and the republic of Venice.”

We are surely not working for the congregation to acquire that kind of visibility? All congregational organizations already possess that kind of visibility.

If there were nothing else to work for, we might begin by rejecting the expression “inward” and all that it includes, and in a short time the congregations would become so secularized that they would be only social clubs; and many congregations have come there: they have lost even the congregation’s outward form and shape.

But for us who work and toil for the congregation, the question will always be whether the inward and invisible stands in proper proportion to the outward, visible growth and activity.

“We must necessarily know what makes us members of the Church, and that living members of it. If we described the Church merely as an outward society of evil and good, men would not understand that Christ’s kingdom is the righteousness of the heart and the impartation of the Holy Spirit; but they would suppose that it consists merely in an outward observance of certain services and customs.” Apology of the Augsburg Confession, Art. 4.

This was said about the Church; but it applies also to the congregation. For everything the Church has, the congregation also has, and everything the Church is, the congregation also is.

It has been, and still remains, the great danger for all congregations and communions to settle down contented with mere outward appearance and condition.

Outwardness was Israel’s great misfortune. They drew near to God with their mouth, but their heart was far away.

We know that the congregation is to appear visibly in the world; but is there room for the Spirit? Is there recognition of sin and repentance in the hearts? Is there a painful sense of our own wretchedness?

Is there no danger of becoming so outwardly directed that, in the midst of all our striving with schools, missions, congregational work, and constitutional teaching, we become stuck fast in what is outward and visible?

Or what is it that is to become visible, and who is it that is to see it?

Scripture says: “The tree is known by its fruit. The fruit of the Spirit is love, joy, peace, patience, kindness, goodness, faithfulness, gentleness, self-control.”

These visible fruits stand in opposition to the works of the flesh and grow forth from the Spirit as the invisible power.

But not everything is genuine that gives itself out to be so. Much is learned and imitated, so that it looks attractive and at times deceptively similar. Through training great results may be attained, so that hypocrites can surpass Christians both in zeal and in self-sacrifice; yes, they can “travel over land and sea to make one adherent.”

Such a thing is visible enough.

But Jesus says:

“Woe to you, scribes and Pharisees, hypocrites! For you are like whitewashed graves, which outwardly appear beautiful, but inwardly are full of dead bones and all uncleanness. So also you outwardly appear righteous to men, but inwardly you are full of hypocrisy and unrighteousness.”

What, then, is it that is to become visible?

It is the invisible that is to become visible.

“But it is by faith, not by thinking, that we understand these things. By faith we understand that the world came into being by God’s Word, so that what is seen did not arise from what is visible.”

The more spiritual a congregation becomes, the more use it has for the Word and the sacraments. Several of us have witnessed this. When faith is strengthened and hearts are kindled through the proper use of the Word and the sacraments, then the congregation receives power and boldness to appear visibly in the world and in word and deed confess Christ.

Then one may well say: “The congregation is visible because God’s Word and the sacraments are there.”

What else should give it the power to become visible?

But Evjen does not like the expressions. They are “some of the most unfortunate,” and so forth.

What, then, is to be done to avert the misfortune and put an end to the confusion?

If the expressions are unbiblical, then we must cease using them about the congregation.

But first another question:

Does the Third Article apply to the congregation?

I hope Evjen will concede that it has very much to do with the congregation, since every Christian congregation has the right to say: “It is a part of my confession of faith.”

There we confess that we believe, and that brings us at once into a field where sight does not reach. For there is that about the congregation which cannot be seen but must be believed. It is the inward and invisible, the work of the Holy Spirit in the hearts, without which the most perfect form and the most ornamented life are only appearance without power, form without content.

On the other hand, I have known congregations and individual persons whose lives were so infected with frailties and defects that one might well be tempted to prefer those smooth, polished people of the world to the rough-grained, sharp-edged Christians who nevertheless were under the discipline and leading of God’s Spirit.

Therefore so much depends upon one’s sight, whether it is right, when spiritual things are to be judged.

Besides, every Christian knows from experience how difficult it is to make the inward and the outward agree.

“What I will, that I do not do; but what I hate, that I do.” Romans 7:15 and following.

This is a long chapter in the history of the congregation and of the individual Christian. But through the Holy Spirit’s testimony concerning what Jesus has done for a sinner, the heart receives peace and joy.

In the Third Article

the congregation confesses:

“I believe in the Holy Spirit, a holy...” and so forth.

God’s Church on earth is one through the one Spirit who dwells and works in all its members. This unity is as real as it is invisible.

What do we see? Church bodies, sects, and parties without end.

But the unity is there according to God’s Word, our faith, and our confession; not visible but invisible, and it is greater and more glorious than we imagine, because it has been brought about by the Holy Spirit.

What has here been said concerning the Church applies also to the congregation insofar as it is a part of the whole; for there too there may be disagreement.

God’s Church is holy; but it is not the holiness of the members that makes it holy, so that its holiness can be seen by all people.

“Why is the Church called holy?

Because the Holy Spirit dwells there and carries out his work of sanctification in all its members; therefore the Church is called holy, despite the sin and frailty found in it.”

The Holy Spirit dwells there; therefore the holiness is far greater and more glorious than we can see and comprehend. But precisely for this reason we confess that we believe it.

Then the invisible and essential thing about the congregation must therefore become an object of faith.

If, then, there were no other and greater unity and holiness in the congregation than that which appears to sight, we would have to despise God’s congregation as the world generally does; for the world does not understand that the essential thing about the congregation is the inward and invisible, which the Holy Spirit works in the hearts of believers.

Believers have the forgiveness of sins and the standing of children before God, and this is no small thing. It is their greatest possession.

But it cannot be seen, at least not at the beginning, because faith is the work of God’s Spirit in the heart, hidden from all men. Yet what is hidden shall be revealed.

But this is to be holy: that God regards us in Christ as though we had never sinned, and this despite our own feeling of sin and guilt and despite the fact that others see our frailty; “for no one is a saint in the eyes of his spouse.”

Believers possess hope.

“For in hope we are saved; but hope that is seen is not hope. Why should anyone hope for what he sees? But if we hope for what we do not see, then we wait for it with patience.”

Believers possess the daily renewal, the mystery of Christianity.

Paul writes to the congregation in Corinth:

“Though our outward man is perishing, yet the inward is renewed day by day.”

Therefore they are not to look at what is visible and lose courage, but to keep the invisible before their eyes and preserve their boldness.

And he writes to the congregation in Colossae, a congregation praised for its faith in Jesus and its love for all the saints:

“You are dead, and your life is hidden with Christ in God.”

“This their life, which is lived in inward fellowship with God, is hidden from the world as Christ himself is.”

In faith and hope, then, the congregation possesses that which is far greater and more glorious than what it can reveal to the world.

“When Christ our life is revealed, then you also shall be revealed with him in glory.”

From Scripture and confession we therefore see that the expressions “visible and invisible, outward and inward” are not only biblical, but indispensable for expressing the inward life and the outward manifestation both in individuals and in congregations.

If, then, we agree that the congregation is both invisible and visible, there is nothing to quarrel about in that matter except what is to receive the greater weight; and that must depend upon whether the one or the other is being passed over in silence or denied.

Life strives toward the light, and true faith must show itself active in love.

If Evjen prefers “an outward side and an inward side,” he is entirely free to do so; but it nevertheless becomes a visible and an invisible side. And thereby no misfortune has occurred.

But it would be a misfortune if the Free Church drifted into formalism by staring itself blind upon the first point concerning “form” and then denying that the congregation is invisible, that is, that it must have God’s Spirit and God’s life in the hearts; for this it must have if there is to be any question at all of a visible congregation.

Therefore this becomes the first and the most important thing.

For as joyful as it is that the congregation grows outwardly, obtains a fine and convenient church building, an able pastor, a well-trained choir, and possesses abundant means so that it can easily carry on its own work and give richly to schools, missions, and works of mercy, all this—as visible as it is—is nevertheless no infallible testimony that matters are equally well with what is inward and invisible.

The Catholics can surpass us in all this and nevertheless stand as a warning sign for God’s congregation against outwardness and formalism.

Concerning this matter Sverdrup wrote:

“It is not our meaning that the congregation’s being built up into God’s temple is an outward constitution whereby the Church is transformed into a well-ordered collection of men who bow to certain outward church laws, and through a definite outward order that sets each in his place from the highest to the lowest. For with all such outward order one may indeed bring into being a splendid building; but it lacks—everything; for it lacks life. God does not dwell therein. It is then not God’s Spirit that joins the whole together, but outward commandments and customs.” Collected Writings II, 97.

But concerning true visibility he says:

“The work of the Holy Spirit shall be known and revealed through the testimony borne by believers. Where God’s children gather for prayer and thanksgiving, for the preaching and hearing of the Word, for the use of the sacraments; where they are active in the Christian home, in the Christian upbringing and religious instruction of the children; where they labor in mission and in the work of mercy; where they speak a word of encouragement to a fellow traveler; where they extend a helping hand in Jesus’ name and for Jesus’ sake—there the Church becomes visible, that is, there its inward invisible being reveals itself, so that it is known that here the mind and Spirit of Christ are present.” Collected Writings II, 235.

The congregation walks by faith, not by sight, and through faith and hope it possesses such infinite riches that neither can it itself see and comprehend them, nor can they be seen and comprehended by the world that surrounds it.

As the body of Christ, it ought to become Christian, that is, like Christ; but this succeeds only in part, for he was perfect.

It is a great thing when we succeed in revealing something of the mind of Christ, in bearing fruit to the Father’s glorification.

But this does not happen through the Law’s: “You shall be visible,” but through faith’s union with Jesus Christ.
